点分树保存点分治的递归关系，每个重心的父亲是上一层重心。因此，点分树中以 $p$ 为根的子树，恰好包含选出 $p$ 时所在连通块的全部点。以下距离均指原树距离。

记点 $y$ 的权值为 $a_y$，查询点 $x$ 时要求 $\sum_y a_y\operatorname{dist}(x,y)$。\texttt{path(x, f)} 从点分树根到 $x$ 枚举祖先 $p$，给出 $d=\operatorname{dist}(x,p)$。当 $p\ne x$ 时，$q$ 是从 $p$ 走向 $x$ 的下一个点。删除重心 $p$ 后，包含 $x$ 的连通块会以 $q$ 为下一层重心，下面称它为分支 $q$；当 $p=x$ 时，$q=0$。

三个数组维护的内容如下：
\begin{itemize}
  \item \texttt{cnt[p]}：选出 $p$ 时，整个连通块内的点权总和。
  \item \texttt{sum[p]}：该连通块内所有 $a_y\operatorname{dist}(p,y)$ 的和。
  \item \texttt{ds[q]}：分支 $q$ 内所有 $a_y\operatorname{dist}(p,y)$ 的和。
\end{itemize}
分支 $q$ 内的点恰好构成点分树中以 $q$ 为根的子树，因此其点权总和就是 \texttt{cnt[q]}。注意，\texttt{ds[q]} 虽然用 $q$ 作下标，存的却是到父重心 $p$ 的加权距离和。

修改点 $x$ 的权值时，只有它的点分树祖先所对应的连通块包含 $x$。沿 \texttt{path} 更新这些祖先的 \texttt{cnt}、\texttt{sum}，并在 $q\ne 0$ 时更新对应分支的 \texttt{ds} 即可。

查询时，在祖先 $p$ 处先加上
\[
\texttt{sum[p]}+\texttt{cnt[p]}\cdot d.
\]
这相当于对当前连通块内每个点 $y$ 加上 $a_y(\operatorname{dist}(p,y)+\operatorname{dist}(x,p))$。若 $p$ 是其中一个端点，或 $x,y$ 位于不同分支，原树路径经过 $p$，这就是正确贡献。

若 $y$ 与 $x$ 同在分支 $q$，路径未必经过 $p$，这一层不统计它。因此，$q\ne 0$ 时减去
\[
\texttt{ds[q]}+\texttt{cnt[q]}\cdot d,
\]
把同一分支刚才加上的部分全部扣掉，留到下一层处理。任意点 $y$ 都在它与 $x$ 首次分属不同分支，或其中一个端点成为当前重心时计入。在更早的层中，它的贡献被扣掉；在更晚的层中，它已不在当前连通块内，因此恰好统计一次。$y=x$ 时在 $p=x$ 处计入，距离贡献为零。

例如原树是链 $1-2-3-4-5-6-7$，第一层选重心 $4$，左右两块分别选 $2$、$6$。查询 $x=3$ 时，各层最终统计的点如下：
\begin{center}
\begin{tabular}{c|c|c}
当前重心 $p$ & 扣掉的分支 & 实际统计的点 \\
\hline
$4$ & $\{1,2,3\}$ & $\{4,5,6,7\}$ \\
$2$ & $\{3\}$ & $\{1,2\}$ \\
$3$ & 无 & $\{3\}$
\end{tabular}
\end{center}
这三组点互不重叠，合起来恰好是全部点。
